\documentclass[10pt,a4paper]{article}
\usepackage[margin=1in]{geometry}
\usepackage{hyperref}
\hypersetup{
    colorlinks=true,
    linkcolor=blue,
    urlcolor=blue,
    citecolor=blue
}
\usepackage{enumitem}
\usepackage{titlesec}
\usepackage{parskip}
\usepackage{fontawesome5}  

\titleformat{\section}{\large\bfseries}{}{0em}{}[\titlerule]

\begin{document}

\begin{center}
    {\Huge \textbf{Fernando Gutiérrez Canales}}\\
    {\Large Desarrollador Python | Ingeniero Backend}\\    
    León, México \quad | \quad \href{https://www.linkedin.com/in/fernando-guti\%C3\%A9rrez-canales-804876343/}{LinkedIn} \quad | \quad +52 1 477 55 18 832 \quad | \quad \href{mailto:carl.cfgc@gmail.com}{carl.cfgc@gmail.com} \\
     \href{https://github.com/Fernando-Canales}{GitHub: github.com/Fernando-Canales}
\end{center}

\vspace{1em}

\section*{Resumen Profesional}
Desarrollador Python con Doctorado en Ciencias Exactas y más de 5 años de experiencia diseñando sistemas \textit{backend} escalables, flujos de trabajo automatizados y \textit{pipelines} de datos robustos. Altamente competente en la integración de Python con C y Bash para optimizar el rendimiento algorítmico y reducir la carga computacional. Sólida base en principios de ingeniería de software, con experiencia práctica utilizando control de versiones (Git), contenedorización (Docker) y bases de datos relacionales (SQL) para entregar código eficiente y listo para producción.

\section*{Experiencia Profesional}
\textbf{Mercor} \hfill \textbf{Göttingen (Remoto)} \\
\textit{Ingeniero de Prompts de IA (Lógica y Pruebas de LLMs)} \hfill \textit{Oct 2025 -- Feb 2026}
\begin{itemize}[leftmargin=1.5em]
	\item Desarrollé conjuntos de datos lógicos y \textit{prompts} adversarios complejos para probar, depurar y refinar (\textit{fine-tune}) Modelos de Lenguaje Grande (LLMs), identificando sistemáticamente casos límite (\textit{edge cases}) y mejorando la capacidad de razonamiento del modelo en un 30\%.
\end{itemize}

\textbf{Observatorio de París \& Instituto Max Planck de Investigaciones del Sistema Solar} \hfill \textbf{París y Göttingen} \\
\textit{Desarrollador Python e Investigador} \hfill \textit{Mar 2022 -- Jul 2025}
\begin{itemize}[leftmargin=1.5em]
    \item Diseñé y mantuve \textit{pipelines} automatizados y escalables en Python para procesar, limpiar y estructurar volúmenes masivos de datos sin procesar para la misión espacial PLATO.
    \item Integré módulos de C y Bash en bases de código Python existentes para optimizar cuellos de botella en el rendimiento, reduciendo drásticamente el tiempo de ejecución y el uso de recursos computacionales.
    \item Colaboré con equipos internacionales en más de 15 instituciones, asegurando la calidad del código, estandarizando arquitecturas de datos y utilizando control de versiones para una integración de proyectos fluida.
\end{itemize}

\textbf{ESTEC, Agencia Espacial Europea (ESA)} \hfill \textbf{Noordwijk, Países Bajos}\\
\textit{Desarrollador Python (Prácticas)}  \hfill \textit{Jun 2023 -- Ago 2023}
\begin{itemize}[leftmargin=1.5em]
    \item Desarrollé e implementé \textit{scripts} automatizados en Python para extraer, validar y analizar parámetros complejos de sensores de hardware, mejorando significativamente la confiabilidad de los flujos de trabajo operativos.
\end{itemize}

\section*{Proyectos Destacados de Desarrollo}
\textbf{Plato's Closet (Paquete de Python Open-Source)} \hfill \textbf{Python, GitHub Actions, Sphinx} \\
\textit{Desarrollador Principal y Mantenedor} \hfill \textit{\href{https://github.com/Fernando-Canales/platos_closet}{github.com/Fernando-Canales/platos\_closet}}
\begin{itemize}[leftmargin=1.5em]
    \item Desarrollé y publiqué un paquete de Python modular e instalable vía \texttt{pip} utilizando \texttt{setuptools}, estableciendo una arquitectura robusta y documentación técnica exhaustiva generada a través de Sphinx.
    \item Implementé \textit{pipelines} automatizados de Integración Continua / Despliegue Continuo (CI/CD) utilizando GitHub Actions para asegurar la calidad del código, ejecutar pruebas automatizadas y optimizar los flujos de lanzamiento (\textit{releases}).
\end{itemize}

\textbf{Arquitectura de Base de Datos para Clínica Oftalmológica} \hfill \textbf{León, México} \\
\textit{Desarrollador Principal (Proyecto Independiente)} \hfill \textit{2021 -- 2022}
\begin{itemize}[leftmargin=1.5em]
    \item Diseñé e implementé una base de datos SQL relacional y centralizada para reemplazar registros físicos y digitales fragmentados, estableciendo una arquitectura \textit{backend} robusta para la clínica.
    \item Desarrollé \textit{scripts} automatizados en Python para consultar la base de datos y generar reportes operativos recurrentes, reemplazando la entrada manual de datos con soluciones programáticas.
\end{itemize}

\section*{Habilidades Técnicas}
\textbf{Lenguajes de Programación:} Python, C, Bash, SQL. \\
\textbf{Ingeniería de Software y Herramientas:} Git, Docker, Linux/Unix, CI/CD (GitHub Actions), Programación Orientada a Objetos (POO), Integración de APIs, \textit{Scripting} y Automatización. \\
\textbf{Librerías y Frameworks:} Pytest, Pandas, NumPy, SciPy, Scikit-learn.

\section*{Educación}
\textbf{Certificación de Bootcamp en Análisis de Datos} \hfill \textit{TripleTen} \hfill \textit{2026}\\
\textbf{Doctorado en Astrofísica} \hfill \textit{Observatorio de París \& MPS Göttingen} \hfill \textit{2025}\\
\textbf{Maestría en Astrofísica} \hfill \textit{Universidad de Guanajuato, México} \hfill \textit{2021}\\
\textbf{Licenciatura en Física} \hfill \textit{Universidad de Guanajuato, México} \hfill \textit{2019}

\section*{Logros y Publicaciones}
\begin{itemize}[leftmargin=1.5em]
    \item Coautor de \href{https://ui.adsabs.harvard.edu/search/fq...}{publicaciones revisadas por pares} enfocadas en análisis computacional complejo y procesamiento automatizado de datos.
    \item Beneficiario de becas académicas completas y graduado con Mención Honorífica en los tres grados universitarios (2019, 2021, 2025).
\end{itemize}

\end{document}
